% paper/wisa2026-zh/main.tex
% WISA 2026 中文短文稿（≤ 7 页）
% 本稿以 article + ctex + twocolumn 近似《计算机工程与应用》投稿排版，
% 录用后须替换为 CEAJ 官方模板。双盲：不含作者、单位、致谢、资助、
% 自引及代码仓链接。
\documentclass[10pt,a4paper,twocolumn]{article}

\usepackage[UTF8,fontset=none]{ctex}
\setCJKmainfont{Noto Serif CJK SC}[
  BoldFont = Noto Serif CJK SC Bold,
  ItalicFont = Noto Serif CJK SC,
  BoldItalicFont = Noto Serif CJK SC Bold
]
\setCJKsansfont{Noto Sans CJK SC}
\setCJKmonofont{Noto Sans Mono CJK SC}
\newCJKfontfamily\heiti{Noto Sans CJK SC Bold}

\usepackage[a4paper,left=2cm,right=2cm,top=2.2cm,bottom=2.2cm,columnsep=0.7cm]{geometry}
\usepackage{amsmath,amssymb}
\usepackage{booktabs}
\usepackage{multirow}
\usepackage[breaklinks=true,hidelinks]{hyperref}
\hypersetup{
  pdftitle={},
  pdfauthor={},
  pdfsubject={},
  pdfkeywords={},
  pdfcreator={},
  pdfproducer={}
}
\usepackage{enumitem}
\usepackage{placeins}
\usepackage{tikz}
\usetikzlibrary{positioning,arrows.meta,calc,patterns}
\usepackage{pgfplots}
\pgfplotsset{compat=1.18}
\setlist[itemize]{topsep=1pt,partopsep=0pt,parsep=1pt,itemsep=1pt,leftmargin=1.6em}
\setlist[enumerate]{topsep=1pt,partopsep=0pt,parsep=1pt,itemsep=1pt,leftmargin=1.8em}
\setlength{\parskip}{2pt plus 0.5pt minus 0.5pt}
\setlength{\textfloatsep}{5pt plus 1pt minus 1pt}
\setlength{\intextsep}{5pt plus 1pt minus 1pt}
\setlength{\abovecaptionskip}{3pt}
\setlength{\belowcaptionskip}{2pt}

\usepackage{titlesec}
\titleformat*{\section}{\normalsize\heiti}
\titleformat*{\subsection}{\small\heiti}
\titlespacing*{\section}{0pt}{1.2ex plus 0.3ex minus 0.2ex}{0.6ex plus 0.1ex}
\titlespacing*{\subsection}{0pt}{0.9ex plus 0.2ex minus 0.2ex}{0.4ex plus 0.1ex}

\providecommand{\doi}[1]{}

\title{\bfseries 面向区块链信息系统执行层的可验证窥孔优化：\\dMIR 进位链算子扩展与 SMT 等价性校验}
\author{匿名作者~~匿名单位}
\date{}

\begin{document}

\twocolumn[
\begin{@twocolumnfalse}
\maketitle
\begin{center}
\parbox{0.9\textwidth}{\small
\textbf{摘要：}区块链 JIT 窥孔优化须逐位保持交易语义，未经等价性
验证的重写规则可能引发共识分叉。LLVM IR 等通用中间表示以多返回值
承载进位、溢出等位级副产物，单条重写规则的前后件跨越多个 IR 结点，
在 Alive/Alive2 中难以表述为单条可验证规则，致使工程侧陷入放弃规则、
未验证部署或人工证明的三难。为确定性虚拟机 DTVM (Deterministic
Virtual Machine) 的中间表示 dMIR (Deterministic Middle IR) 提出
IR 侧一等位向量算子扩展：在不改求解器前提下，将位级运算模式建模为
dMIR 一等算子，纳入 Z3 位向量等价性判据。基于此实现两级可验证
窥孔重写：dMIR 层 \textbf{70} 条规则经 Z3 逐条等价性校验，x86 层
\textbf{13} 条机械冗余清理规则以对照测试套件覆盖，并构建枚举
\textbf{853} 条候选的合成管线。evmone-bench \textbf{27} 项基准几何
平均加速 \textbf{7.24\%}，峰值 \textbf{25.79\%}；\textbf{5884} 条
Cancun 差分测试中启用/关闭两侧的账户状态、存储与 gas 消耗逐字节
一致。
\vspace{0.3em}

\textbf{关键词：}区块链信息系统；EVM JIT；dMIR；窥孔优化；SMT 等价性校验；进位链
}
\end{center}
\vspace{1.0em}
\end{@twocolumnfalse}
]



\end{document}
